\chapter{Groupoids in algebraic spaces}

In this chapter we introduce groupoids in algebraic spaces, which provide the
correct framework to describe quotients of algebraic stacks and play a central
role in the proof of the Keel--Mori theorem.

\section{Groupoids in Algebraic Spaces}

\begin{definition}[Groupoid in algebraic spaces]
	(\cite[][Tag 043V]{stacks-project})
	Let $B$ be a base algebraic space. A \emph{groupoid in algebraic spaces over $B$}
	consists of a quintuple
	\[
	(U, R, s, t, c)
	\]
	where:
	\begin{itemize}
		\item $U$ and $R$ are algebraic spaces over $B$,
		\item $s,t : R \to U$ are morphisms (source and target),
		\item $c : R \times_{s,U,t} R \to R$ is a composition morphism,
	\end{itemize}
	such that for every scheme $T \to B$, the induced data
	\[
	(U(T), R(T), s, t, c)
	\]
	define a groupoid (i.e. a category in which all morphisms are invertible).
\end{definition}

\begin{remark}
	The definition is functorial in $T$. Hence all axioms of a groupoid can be
	checked on $T$-valued points.
\end{remark}

\begin{lemma}[Identity and inverse]
	(\cite[][Tag 043V]{stacks-project})
	Let $(U,R,s,t,c)$ be a groupoid in algebraic spaces over $B$. Then there exist
	unique morphisms
	\[
	e : U \to R, \qquad i : R \to R
	\]
	such that for every scheme $T \to B$:
	\begin{itemize}
		\item $e$ defines identity morphisms,
		\item $i$ defines inverses.
	\end{itemize}
	
	We will write $(U,R,s,t,c,e,i)$ for the corresponding septuple.
	We write
	\[
	R \rightrightarrows U
	\]
	for a groupoid in algebraic spaces, where the two arrows are given by
	$s,t : R \to U$.
\end{lemma}


\section{Relation to Equivalence Relations}

\begin{definition}[Associated relation]
	Let $(U,R,s,t,c)$ be a groupoid in algebraic spaces over $B$. Then the morphism
	\[
	j = (t,s) : R \to U \times_B U
	\]
	is called the associated \emph{pre-equivalence relation}.
\end{definition}

\begin{lemma}
	(\cite[][Tag 043V]{stacks-project})
	The morphism
	\[
	j : R \to U \times_B U
	\]
	associated to a groupoid defines a pre-equivalence relation on $U$.
\end{lemma}

\begin{remark}
	Thus groupoids generalize equivalence relations:
	\begin{itemize}
		\item If $j$ is a monomorphism, one recovers an equivalence relation.
		\item In general, groupoids allow nontrivial automorphisms (stabilizers).
	\end{itemize}
\end{remark}

---

\section{Examples}

\begin{example}[Equivalence relations]
	Let $j : R \to U \times_B U$ be an equivalence relation in algebraic spaces.
	Then there is a unique groupoid structure on $(U,R)$ extending this data.
\end{example}

\begin{example}[Group actions]
	Let $G$ be a group algebraic space over $B$ acting on an algebraic space $U$.
	Then
	\[
	R := G \times_B U
	\]
	with
	\[
	s(g,x) = x, \quad t(g,x) = g \cdot x
	\]
	defines a groupoid in algebraic spaces.
\end{example}

---

\section{Stabilizers}

\begin{definition}[Stabilizer]
	(\cite[][Tag 0448]{stacks-project})
	Let $(U,R,s,t,c)$ be a groupoid in algebraic spaces over $B$. The
	\emph{stabilizer} is the algebraic space
	\[
	G := R \times_{(t,s), U \times_B U, \Delta} U,
	\]
	where $\Delta : U \to U \times_B U$ is the diagonal.
\end{definition}

\begin{remark}
	The stabilizer parametrizes automorphisms of objects. It is the groupoid-theoretic
	counterpart of the inertia stack of a stack.
\end{remark}

---

\section{Base Change}

\begin{lemma}[Base change of groupoids]
	(\cite[][Tag 043V]{stacks-project})
	Let $(U,R,s,t,c)$ be a groupoid over $B$, and let $B' \to B$ be a morphism of
	algebraic spaces. Then
	\[
	U' := B' \times_B U, \qquad R' := B' \times_B R
	\]
	define a groupoid $(U',R',s',t',c')$ over $B'$.
\end{lemma}

---

\section{Quotients (Preview)}

\begin{definition}[Sheaf quotient]
	Let $(U,R)$ be a groupoid in algebraic spaces over $B$. The presheaf quotient
	is defined by
	\[
	(U/R)(T) := U(T)/R(T).
	\]
	Its sheafification (in the étale topology) is called the \emph{quotient sheaf}.
\end{definition}

\begin{remark}
	A central problem is to determine when $U/R$ is representable by an algebraic
	space. This will be the key step in the proof of the Keel--Mori theorem.
\end{remark}

---

\section{Relation to Stacks}

\begin{remark}[Groupoids and stacks]
	Let $\mathcal{X}$ be an algebraic stack and let $U \to \mathcal{X}$ be a smooth
	surjective morphism. Then
	\[
	R := U \times_{\mathcal{X}} U
	\]
	carries a natural structure of a groupoid in algebraic spaces:
	\[
	R \rightrightarrows U.
	\]
	
	This groupoid presents the stack $\mathcal{X}$ as a quotient:
	\[
	\mathcal{X} \simeq [U/R].
	\]
\end{remark}